\documentclass[12pt]{article}
\usepackage[utf8]{inputenc}

% Load one of the two bibliography style package with the 'author-year' or the 'sequential-numbered' option
% Load one of the two geometry style packages with the 'manuscript' (pre-accetance) or the 'publication' (post-acceptance) option
\usepackage[manuscript,author-year]{fqxi_foundational_review}
\setcitestyle{aysep={}}

\newcommand{\fqxi}{FQ$_x$I}

\title{FQ$_x$I Template for submissions to the Foundational Review Series}

\author[1,\orcidlink{0000-0002-9626-7257}]{Nathan Rutherford \email{nathan.rutherford@fqxi.org}}
\author[1,2, \orcidlink{0000-0003-0695-8220}]{Fabio Anza \email{fabio.anza@fqxi.org}}


\affil[1]{Foundational Questions Institute}
\affil[2]{University of Maryland, Baltimore County}

\date{}

\begin{document}

\maketitle

% TODO: Find a way to create a table of contents. 
\begin{Abstract}
This is where the abstract goes    
\end{Abstract}

\newpage

\tableofcontents

\newpage

\section{Introduction}\label{sec:intro}
The citation commands are \cite{Rutherford24} or \citep[][]{Rutherford24}. You can also add text before and after a citation if you wish \citep[see][for further details]{Rutherford24}.



Here we have two equations

\begin{equation}\label{eq:1}
    R_{\mu \nu} - \frac{1}{2}g_{\mu \nu} R = \frac{8 \pi G}{c^4} T_{\mu \nu}.
\end{equation}

\begin{equation}\label{eq:2}
    E = mc^2
\end{equation}

\section{This is another section}

\subsection{a sub-section}
In Sec.~\ref{sec:intro} we introduced the topic. In Eq.~(\ref{eq:1}) and (\ref{eq:2}), we show Einstein's field equations and rest mass-energy equivalence relation, respectively.

\bottomline{Give a brief 1-3 sentence summary of what a reader should be taking away from this section. }

\section*{Acknowledgments}
Give thanks to anyone involved in helping with the manuscript, \fqxi, and any other funding sources. 
\appendix

\section{Appendix A}

\setfqxibib
\bibliography{biblio.bib}
\end{document}
